
\medskip

\begin{enumerate}
\item 
	\begin{enumerate}
		\item %Justifier que $330$ n'est pas un nombre premier.
		330 est pair : il n'est donc pas premier (le seul premier pair est 2).
		
%La décomposition en produit de facteurs premiers de $500$ est: $500 = 2^2 \times 5^3$.
		\item %Décomposer $330$ en produit de facteurs premiers.
$330 = 10 \times 33 = 2 \times 5 \times 3 \times 11 = 2 \times 3 \times 5 \times 11$.
		\item %Justifier que $165$ divise $330$.
$330 = 2 \times 165$, donc 165 est un diviseur de 330.
		\item Justifier que 165 ne divise pas $500$.
$165 = 15 \times 11 = 3 \times 5 \times 11$, donc 11 divise 165, mais 11 n'est pas un diviseur de 500 (11 n'est pas dans la liste des diviseurs premiers de 500).
	\end{enumerate}
\end{enumerate}
	
%La pâtisserie Délices a préparé $330$ biscuits aux noix et $500$ biscuits au chocolat.

%\textbf{La pâtisserie souhaite répartir le plus de biscuits possible dans $165$ boites.}

%La pâtisserie met le même nombre de biscuits aux noix dans chaque boîte.

\begin{enumerate}[resume]
\item %Combien de biscuits aux noix y a-t-il dans chaque boîte ?
On a $330 = 165 \times 2$ : on peut donc mettre 2 biscuits aux noix dans chacune des 165 boîtes.

%La pâtisserie met aussi le même nombre de biscuits au chocolat dans chaque boîte.
\item
	\begin{enumerate}
		\item %Combien de biscuits au chocolat y a-t-il dans chaque boîte ?
		On a $500 = 165 \times 3 + 5$ : on peut donc mettre 3 biscuits au chocolat dans chaque boîte.
		\item Combien de biscuits au chocolat reste-t-il ?
	\end{enumerate}
		
%Une boîte de biscuits coûte \np{3650} francs.

%À partir de 10 boîtes achetées, la pâtisserie Délices offre une réduction de 5\,\% sur le montant total.
\item %Combien va-t-on payer pour l'achat de 12 boîtes ?

%Faire apparaÎtre les calculs effectués.
Retrancher 5\,\% c'est multiplier par $1 - \dfrac{5}{100} = 1 - 0,05 = 0,95$.

À partir de 10 boîtes achetées chaque boîte est donc facturée $\np{3650} \times 0,95$.

Pour 12 boîtes achetées le prix effectivement payé sera :

$12 \times \np{3650} \times 0,95 = \np{42610}$~(\euro).

\end{enumerate}

\bigskip

